% ============================================================
%  Глава 11 — 3D-пикер цели: математика угла
% ============================================================
\chapter{3D-пикер цели: математика}
\label{ch:picker}

\section{Вопрос}

Пользователь хочет указать курс $\varphi$, на который робот развернётся
перед движением. Как именно? Можно было бы поставить поле «введите угол
в градусах», но это \emph{неудобно} --- человек редко интуитивно знает,
как $+72^\circ$ соотносится с «вон туда, под небольшим углом влево».

\begin{ideabox}
\textbf{3D-пикер} --- решение из feat/mps-target-picker. На экране
показана 3D-сцена: робот в центре, окружность фиксированного радиуса
вокруг него, маркер на окружности. Пользователь \emph{кликает} в любую
точку «пола» сцены --- маркер переезжает на окружность в направлении
клика, а угол $\varphi$ автоматически вычисляется.

Это даёт мгновенную визуальную обратную связь и интуитивный ввод.
Под капотом --- немного линейной алгебры и одна функция \texttt{atan2}.
\end{ideabox}

\section{Геометрия и соглашения}

В 3D-сцене Three.js (React Three Fiber) система координат другая, чем
в мире робота. Аккуратно зафиксируем:

\begin{center}
\begin{tabular}{lll}
\toprule
& \textbf{Мир робота} & \textbf{Three.js}\\
\midrule
$x$ ось & вперёд от робота (curr. heading)  & вправо от камеры\\
$y$ ось & влево от робота & вверх (высота над полом)\\
$z$ ось & вверх & к камере (от сцены к зрителю)\\
\bottomrule
\end{tabular}
\end{center}

Робот «смотрит» вдоль своей $+x$ при yaw $= 0$. В Three угол сцены
сделан так, что пол --- плоскость $y = 0$ Three; камера слегка выше
смотрит вниз.

\textbf{Соответствие координат:} мир $(w_x, w_y)$ соответствует Three
$(w_x, h, -w_y)$, где $h$ --- высота над полом (для маркера ---
константа \texttt{MARKER\_HEIGHT}). Минус в третьей координате потому
что Three-$z$ направлен «к камере», а мировой $y$ --- «влево».

\begin{notebox}
Соглашение об углах: $\varphi$ --- относительный курс. $\varphi = 0$
--- прямо вперёд (вдоль текущего yaw робота); $\varphi > 0$ --- влево
(CCW против часовой); $\varphi < 0$ --- вправо. Это совпадает с
соглашением ROS и нашим robot frame в коде.
\end{notebox}

\section{Точка клика \texorpdfstring{$\to$}{->} угол}

\subsection{Формула}

В Three.js «raycast» от позиции мыши даёт точку клика на полу
$(x_3, z_3)$. Переводим в мировые $(w_x, w_y) = (x_3, -z_3)$. Угол ---
обычная $\operatorname{atan2}$:

\begin{formulabox}[title={Угол от точки клика}]
\begin{equation}
  \varphi = \operatorname{atan2}(w_y, w_x) = \operatorname{atan2}(-z_3, x_3).
  \label{eq:click_to_angle}
\end{equation}
\end{formulabox}

Функция \texttt{atan2} (есть в любой математической библиотеке)
аккуратно ставит знаки и выдаёт результат в $(-\pi, \pi]$. Заодно
не путается на $w_x = 0$ (где «наивный» $\arctan(w_y/w_x)$ дал бы
деление на ноль).

\begin{codebox}[title={\filepath{compute\_node/frontend/src/lib/targetAngle.ts}}]
\begin{lstlisting}[style=tsstyle]
// Точка клика на полу сцены (Three-координаты x, z) → относительный курс φ.
export function groundPointToAngle(
  threeX: number, threeZ: number
): number {
  // мир: wx = threeX, wy = -threeZ. φ = atan2(wy, wx).
  return Math.atan2(-threeZ, threeX)
}
\end{lstlisting}
\end{codebox}

\subsection{Почему именно \texttt{atan2}?}

\begin{ideabox}
\textbf{Альтернатива --- $\operatorname{atan}(w_y/w_x)$.} Беда:
\begin{itemize}
  \item Делит на ноль при $w_x = 0$ (точно «влево» или «вправо»).
  \item Не различает квадранты: $\arctan(1/1) = \arctan(-1/-1) = \pi/4$,
        хотя углы --- $45^\circ$ и $-135^\circ$.
\end{itemize}

\texttt{atan2} принимает оба аргумента \emph{отдельно}, использует знаки,
покрывает все 4 квадранта и не падает на нуле. Это стандартный приём в
графике и робототехнике.
\end{ideabox}

\subsection{Численные примеры}

\begin{examplebox}
$(w_x, w_y) = (1, 0)$: точка прямо перед роботом. $\varphi = \operatorname{atan2}(0, 1) = 0$ --- «прямо».

$(w_x, w_y) = (0, 1)$: точка слева от робота. $\varphi = \operatorname{atan2}(1, 0) = \pi/2 \approx 90^\circ$.

$(w_x, w_y) = (-1, 0)$: точка прямо за роботом. $\varphi = \operatorname{atan2}(0, -1) = \pi$ --- разворот на $180^\circ$.

$(w_x, w_y) = (-1, -1)$: точка справа-сзади. $\varphi = \operatorname{atan2}(-1, -1) = -3\pi/4 \approx -135^\circ$.
\end{examplebox}

\section{Угол \texorpdfstring{$\to$}{->} позиция маркера}

Обратное: имея $\varphi$ и фиксированный радиус $N$ (радиус окружности
в сцене), найти позицию маркера в Three-координатах.

\begin{formulabox}[title={Маркер на окружности радиуса $N$}]
\begin{equation}
  (x_3, h, z_3) = (N\cos\varphi,\; h_\text{marker},\; -N\sin\varphi).
  \label{eq:angle_to_pos}
\end{equation}
\end{formulabox}

Объяснение: в мировой системе точка на окружности радиуса $N$ под
углом $\varphi$ --- это $(N\cos\varphi, N\sin\varphi)$. Переходим в
Three: $z_3 = -w_y = -N\sin\varphi$, $x_3 = w_x = N\cos\varphi$.
Высота $h_\text{marker}$ --- константа над полом (чтобы маркер не
тонул в пол).

\begin{codebox}[title={\filepath{lib/targetAngle.ts}}]
\begin{lstlisting}[style=tsstyle]
export const MARKER_HEIGHT = 0.03

export function angleToMarkerPosition(
  angle: number,
  radius: number,
): [number, number, number] {
  return [radius * Math.cos(angle), MARKER_HEIGHT, -radius * Math.sin(angle)]
}
\end{lstlisting}
\end{codebox}

\subsection{Почему радиус \emph{фиксированный}?}

В предыдущей версии (одна фаза, без TURN) пикер позволял указать
\emph{любую точку} --- и расстояние, и направление. Это перегружало
интерфейс: пользователь думает «куда» И «насколько далеко» одновременно.

В двухфазной схеме \emph{дистанция $D$ задаётся отдельным полем},
а пикер отвечает только за \emph{направление}. Поэтому маркер всегда
ползёт по окружности фиксированного радиуса $N$ --- независимо от того,
куда конкретно кликнул пользователь. Это \emph{разделение
ответственности}: одна задача --- один контрол.

\section{Центральная мёртвая зона}

\begin{ideabox}
\textbf{Проблема:} если кликнуть \emph{рядом с центром}, маленькое
движение мыши даёт огромный скачок $\varphi$. У точки $(0.01, 0.01)$
угол $= \pi/4$, у точки $(0.01, -0.01)$ угол $= -\pi/4$ --- разница в
$90^\circ$ при микроперемещении в $1$~см.

\textbf{Решение:} в зоне радиуса $r_\text{dead} = 0.3\,N$ от центра
\emph{игнорируем} клик.
\end{ideabox}

\begin{codebox}[title={\filepath{lib/targetAngle.ts}}]
\begin{lstlisting}[style=tsstyle]
export const CENTER_DEADZONE_FRACTION = 0.3

export function groundPointRadius(
  threeX: number, threeZ: number
): number {
  return Math.hypot(threeX, threeZ)
}
\end{lstlisting}
\end{codebox}

В UI-логике: если \texttt{groundPointRadius < CENTER\_DEADZONE\_FRACTION * N},
клик \emph{не} обновляет $\varphi$.

\section{Подпись угла}

Финальный штрих --- человекочитаемая подпись маркера: «прямо», «+35°»,
«-40°»:

\begin{codebox}[title={\filepath{lib/targetAngle.ts}}]
\begin{lstlisting}[style=tsstyle]
export function formatHeadingLabel(angle: number): string {
  const deg = (angle * 180) / Math.PI
  if (Math.abs(deg) < 3) return 'прямо'
  const rounded = Math.round(deg)
  return rounded > 0 ? `+${rounded}°` : `${rounded}°`
}
\end{lstlisting}
\end{codebox}

Дед-зона по \emph{углу} ($\le 3^\circ$ показывается как «прямо») --- не
путать с дед-зоной по \emph{радиусу}. Это разные защиты:

\begin{itemize}
  \item Центральная (радиус) --- против скачков при микроперемещениях.
  \item Угловая (3°) --- косметика лейбла, чтобы не мигало
        $+1^\circ / -1^\circ$ при дрожании руки.
\end{itemize}

\section{Полный пайплайн пикера}

\begin{center}
\begin{tikzpicture}[
  node distance=8mm,
  pipestep/.style={
    draw, rounded corners=2pt, thick, fill=blue!8,
    align=center, font=\footnotesize, minimum width=4.5cm, minimum height=8mm,
  },
  flow/.style={-{Stealth[length=2mm]}, thick},
]
  \node[pipestep] (s1) {\textbf{1.} Клик мыши $(p_x, p_y)$};
  \node[pipestep, below=of s1] (s2) {\textbf{2.} Raycast $\to$ $(x_3, z_3)$ на полу};
  \node[pipestep, below=of s2] (s3) {\textbf{3.} Дед-зона: $\sqrt{x_3^2+z_3^2} < 0.3N$?};
  \node[pipestep, below=of s3] (s4) {\textbf{4.} $\varphi = \operatorname{atan2}(-z_3, x_3)$};
  \node[pipestep, below=of s4] (s5) {\textbf{5.} Маркер $(N\cos\varphi, h, -N\sin\varphi)$};
  \node[pipestep, below=of s5] (s6) {\textbf{6.} Лейбл «прямо» / «+35$^\circ$»};
  \node[pipestep, below=of s6, fill=green!8] (s7) {\textbf{7.} На «Run»: $\varphi$ в \texttt{target\_heading}};
  \draw[flow] (s1) -- (s2);
  \draw[flow] (s2) -- (s3);
  \draw[flow] (s3) -- node[right, font=\tiny]{нет: продолжаем} (s4);
  \draw[flow] (s4) -- (s5);
  \draw[flow] (s5) -- (s6);
  \draw[flow] (s6) -- (s7);
  \draw[flow] (s3.east) .. controls +(2,0) and +(2,-1) ..
              node[right, font=\tiny]{да: игнор} (s1.east);
\end{tikzpicture}
\end{center}

\section{От UI к роботу}

Когда пользователь жмёт «Run on Robot», значение $\varphi$ уходит в
REST-запрос:

\begin{codebox}[title={REST}]
\begin{verbatim}
POST /api/v1/mps/scenario/run
Content-Type: application/json

{
  "source": "robot",
  "request": {
    "distance": 2.0,
    "v_target": 0.15,
    "target_heading": 0.7854   // <-- phi (угол) из пикера, рад
  }
}
\end{verbatim}
\end{codebox}

\texttt{routers/mps.py} публикует это в MQTT \texttt{mps/scenario/run},
\texttt{mps\_node} читает, валидирует ($-\pi \le \varphi \le \pi$),
ставит \texttt{run.target\_heading = $\varphi$}, фаза \texttt{'turn'}
--- и поехали.

\section{Тесты математики пикера}

Файл \filepath{compute\_node/frontend/src/lib/targetAngle.test.ts}
(см. в проекте) покрывает:

\begin{itemize}
  \item Прямые тестовые случаи: $(1, 0) \to 0$, $(0, -1) \to \pi/2$ (помним
        Three's $z$ перевёрнут), $(-1, 0) \to \pi$ и т.д.
  \item Симметрию: \texttt{angleToMarkerPosition(groundPointToAngle(x, z) | radius = N)}
        возвращает точку на окружности радиуса $N$ \emph{в том же
        направлении}, что исходный клик.
  \item Граничные значения подписи: $\pm 2.9^\circ \to$ «прямо»,
        $\pm 3.1^\circ \to$ числовая подпись.
  \item Дед-зону по радиусу.
\end{itemize}

Юнит-тесты в jsdom без WebGL --- именно потому модуль \texttt{targetAngle.ts}
\emph{не зависит} от Three: вся «математика» вынесена в чистые функции,
а 3D-визуал --- отдельным компонентом \texttt{MpsTargetScene.tsx}.

\section{Что мы получили}

\begin{itemize}
  \item Пикер цели --- интуитивный визуальный ввод угла $\varphi$.
  \item Математика тривиальная: $\varphi = \operatorname{atan2}(-z_3, x_3)$
        для клика $\to$ угла, $(N\cos\varphi, h, -N\sin\varphi)$ для угла $\to$ позиции.
  \item Соглашения о координатах фиксированы и совпадают с тем, что
        потом получает Pi.
  \item Дед-зона по радиусу защищает от скачков, дед-зона по углу
        --- косметика.
  \item Юнит-тестируется без WebGL --- спасибо чистым функциям.
\end{itemize}

\medskip
\textbf{Это была вся «активная» математика модуля МПС.} Дальше осталась
ещё одна, инфраструктурная глава --- как матрицы из MATLAB попадают
в Python.
